\section{Prospective Routing-Resolution Study}
\label{sec:routing-resolution}

\subsection{Frozen design}

The follow-up estimand and source were fixed before a fresh 256-bit root seed was
generated.  The freeze, seed generation, and one-shot execution ran in that order
in one scripted sequence; integrity follows from the frozen hashes and seed
commitment, not from elapsed delay.
Each of 120 independent worlds has five coordinates modulo 7 and one graph
motif consisting of two distinct sources entering a common target.  The 30
possible motifs are crossed with edge heads
\[
 a_s(1+x_t),\qquad a_s(1+x_t)^2,\qquad
 a_s(1+x_s+x_t).
\]
One third of worlds use two original linear heads; two thirds use at least one
of the other heads.  Direct multipliers and two edge coefficients are fitted
from 800 primitive-action transitions with coordinate noise 0.08.  Graph and
head family are selected using 400 disjoint transitions.  No fitting reads the
1,200 OOD transitions, all of which have two active action coordinates and
coordinate noise 0.12.  The designated composition stratum activates both
causal edges simultaneously.

The criterion comparison uses the same learned predictor under two routing
maps.  During the first three steps of a six-step stochastic trajectory, the
correct score records target-coordinate discrepancies when either learned
source is active; the negative control routes the same discrepancy field
through a fixed shifted graph.  A logistic map fitted on 24 public development
worlds was frozen before confirmation.  Its outcome is terminal mismatch at
the learned target, and its estimand is wrong-routing minus correct-routing
Brier score.  Thus the outcome does not contain either criterion score, the
diagnostic precedes the terminal outcome, and the two arms differ only in
routing.

Eight inferential quantities use Bonferroni simultaneous two-sided intervals
with familywise level 0.05.  The identification gate additionally requires a
simultaneous Wilson lower bound of 0.90.  Point-effect thresholds are 0.04 for
the learned composition NLL contrast, 0.025 for rollout Hamming area, and 0.005
for criterion Brier score.  These point thresholds are distinct from the null
boundary used by the simultaneous intervals.

The analysis object reports ten interval-bearing quantities at the frozen
Bonferroni divisor of eight.  They contain one exact duplicate pair, one
sign-flipped pair, and one deterministic zero, leaving seven informationally
distinct stochastic quantities.  The frozen v1 contract specified only the
integer divisor and did not enumerate its members; this is a reporting defect,
not a preregistered family definition.  A post-review sensitivity using divisor
ten leaves every decision unchanged.  The cohort, root seed, and multiplicity
family are shared with Paper~4; the three endpoints reported here are a subset
of that single frozen analysis, not independent evidence.

No power analysis for 120 worlds was frozen before confirmation.  A
post-review retrospective design justification, using only the 24 development
worlds and 20,000 stratified bootstrap draws, estimates conjunctive nine-gate
power at 0.9694 with Monte Carlo interval [0.9670, 0.9718].  This calculation
does not use confirmatory results and cannot retroactively establish a
preregistered sample-size rationale.

\subsection{Confirmatory results}

The learned motif and both head families were recovered in all 120 worlds
(rate 1.000; simultaneous Wilson lower bound 0.941).  Table~\ref{tab:resolution}
reports the three Paper~3 endpoints.  Every interval is computed across worlds;
cases and rollouts are nested observations.

\begin{table}[ht]
\centering
\caption{Prospective learned-routing effects.  Positive values favor learned
or correct routing.}
\label{tab:resolution}
\begin{tabular}{lrrr}
\toprule
Endpoint & Mean & World SD & Simultaneous interval \\
\midrule
Composition NLL & 0.35933 & 0.02012 & [0.35431, 0.36435] \\
Rollout Hamming area & 0.08996 & 0.01162 & [0.08706, 0.09286] \\
Routing-correct Brier & 0.02402 & 0.02381 & [0.01808, 0.02997] \\
\bottomrule
\end{tabular}
\end{table}

All three point effects exceed their frozen thresholds and all simultaneous
lower bounds exceed zero.  Performance on the 80 worlds outside the original
linear family was noninferior to the 40 linear worlds: outside-minus-linear NLL
was -0.00132 with simultaneous interval [-0.01031, 0.00766], whose upper bound
is below the 0.03 margin.  This is a direct learned-routing criterion result,
not a reinterpretation of the earlier four-arm average.

In a separate post-review descriptive sensitivity, train noise
\(\{0.04,0.08,0.16\}\) was crossed with OOD noise
\(\{0.06,0.12,0.24\}\), with 120 public-seed worlds per cell.  Graph and head
recovery remained 1.000 in all nine cells.  The minimum mean learned-routing
NLL effect was 0.23337.  This panel was not confirmatory and varies only the
declared coordinate-corruption process.
